% ============================================================
\chapter{$T_0$-Spaces and Specialization Order}
\label{ch:t0_specialization}
% ============================================================

One of the most fruitful discoveries in twentieth-century topology is that
behind every topological space hides an \emph{order}. This order, called the
``specialization order,'' connects topology to partial order theory in a
deep and unexpected way. In the world of $T_0$-spaces, this correspondence
becomes an equivalence: the topology \emph{is} the order, and the order
\emph{is} the topology. It is this duality that makes $T_0$-spaces the
central objects of this course.

\begin{intuition}
The specialization order is the fundamental bridge between topology and order theory. In a
$T_0$-space, each point occupies a position in a partial order determined by the topology.
This chapter systematically explores this correspondence and its consequences.
\end{intuition}

% ------------------------------------------------------------
\section{The Specialization Order}
% ------------------------------------------------------------

\begin{definition}[Specialization order]
Let $(X,\tau)$ be a topological space. The \emph{specialization order} is the relation
$\leq$ on $X$ defined by:
\[
  x \leq y \quad \Longleftrightarrow \quad x \in \overline{\{y\}}
  \quad \Longleftrightarrow \quad \text{every open set containing } x
  \text{ also contains } y.
\]
\end{definition}

\begin{proposition}[Basic properties]
Let $(X,\tau)$ be a topological space with specialization order $\leq$.
\begin{enumerate}
  \item $\leq$ is a preorder (reflexive and transitive).
  \item $\leq$ is a partial order (antisymmetric) if and only if $X$ is $T_0$.
  \item $x \leq y$ if and only if $\overline{\{y\}} \subseteq \overline{\{x\}}$
    (note the reversal!).
  \item Every open set $U$ is upward-closed: if $x \in U$ and $x \leq y$, then $y \in U$.
  \item Every closed set $F$ is downward-closed: if $y \in F$ and $x \leq y$, then
    $x \in F$.
\end{enumerate}
\end{proposition}

\begin{proof}
(1) Reflexivity: $x \in \overline{\{x\}}$ always. Transitivity: if $x \leq y$ and
$y \leq z$, then for every open $U \ni x$, we have $y \in U$ (since $x \leq y$), then
$z \in U$ (since $y \leq z$), so $x \leq z$.

(2) Antisymmetry: $x \leq y$ and $y \leq x$ mean $\overline{\{x\}} = \overline{\{y\}}$,
i.e., $x$ and $y$ are topologically indistinguishable. This equals $x = y$ iff $X$ is
$T_0$.

(4) Let $U$ be open, $x \in U$, $x \leq y$. Every open containing $x$ contains $y$; in
particular $U$ does, so $y \in U$.

(5) By taking complements.
\end{proof}

\begin{warning}
The ordering convention varies among authors. We follow the convention where $x \leq y$
means ``$x$ specializes to $y$,'' i.e., $x \in \overline{\{y\}}$, or equivalently, $y$ is
``more generic'' than $x$. Some authors use the opposite convention. In our convention,
open sets are \emph{upward-closed}.
\end{warning}

% ------------------------------------------------------------
\section{Alexandrov Topology}
% ------------------------------------------------------------

\begin{definition}[Alexandrov topology]
Let $(P, \leq)$ be a preordered set. The \emph{Alexandrov topology} on $P$ is the topology
whose open sets are exactly the upward-closed subsets:
\[
  U \in \tau_{\mathrm{Alex}} \quad \Longleftrightarrow \quad
  (x \in U \text{ and } x \leq y \Rightarrow y \in U).
\]
\end{definition}

\begin{theorem}[Characterization of Alexandrov topology]
A topological space $(X,\tau)$ carries the Alexandrov topology of its specialization order
if and only if $\tau$ is closed under arbitrary intersections.
\end{theorem}

\begin{proof}
$(\Rightarrow)$ Upward-closed subsets of a preorder are closed under arbitrary
intersections.

$(\Leftarrow)$ If $\tau$ is closed under arbitrary intersections, then for every $x$,
$U_x = \bigcap\{U \in \tau : x \in U\}$ is open. This is the smallest open set containing
$x$. We have $y \in U_x$ iff every open containing $x$ contains $y$, i.e., $x \leq y$.
Thus $U_x = \uparrow\! x$ and every open is a union of such sets, hence upward-closed.
\end{proof}

\begin{proposition}[Equivalence of categories]
The functor associating to a preordered set $(P,\leq)$ the Alexandrov space
$(P, \tau_{\mathrm{Alex}})$ establishes an equivalence of categories between
$\mathbf{Preord}$ (preorders and monotone maps) and the full subcategory of
$\mathbf{Top}$ consisting of Alexandrov spaces. Restricting to partial orders gives an
equivalence with $T_0$ Alexandrov spaces.
\end{proposition}

\begin{proof}
A map $f\colon P \to Q$ is monotone if and only if the preimage of every upward-closed set
is upward-closed, which is equivalent to continuity for the Alexandrov topologies. The
inverse functor associates to an Alexandrov space its specialization order.
\end{proof}

% ------------------------------------------------------------
\section{Closure and Order}
% ------------------------------------------------------------

\begin{proposition}[Closure and downward-closed sets]
In any topological space $(X,\tau)$:
\begin{enumerate}
  \item $\overline{\{x\}} = \downarrow\! x$ (for the specialization order).
  \item For any $A \subseteq X$, $\overline{A} = \downarrow\! A$ if and only if $X$
    carries the Alexandrov topology.
  \item In general, $\downarrow\! A \subseteq \overline{A}$, with equality for singletons.
\end{enumerate}
\end{proposition}

\begin{proof}
(1) By definition, $y \in \overline{\{x\}}$ iff every open containing $y$ contains $x$,
i.e., $y \leq x$, i.e., $y \in \downarrow\! x$.

(3) If $a \in A$ and $y \leq a$, then $y \in \overline{\{a\}} \subseteq \overline{A}$,
so $\downarrow\! A \subseteq \overline{A}$.
\end{proof}

\begin{definition}[Generic point]
Let $F$ be an irreducible closed set of $X$ (i.e., $F \neq \emptyset$ and if
$F = F_1 \cup F_2$ with $F_1, F_2$ closed, then $F = F_1$ or $F = F_2$). A point
$x \in F$ is a \emph{generic point} of $F$ if $\overline{\{x\}} = F$.
\end{definition}

\begin{remark}
The existence and uniqueness of generic points for irreducible closed sets is at the heart
of the notion of sober space (Chapter~3).
\end{remark}

% ------------------------------------------------------------
\section{Continuity and Monotonicity}
% ------------------------------------------------------------

\begin{theorem}[Continuity implies monotonicity]
If $f\colon X \to Y$ is continuous, then $f$ is monotone for the specialization orders:
$x \leq_X y \Rightarrow f(x) \leq_Y f(y)$.
\end{theorem}

\begin{proof}
Let $V$ be an open set in $Y$ containing $f(x)$. Then $f^{-1}(V)$ is open in $X$ and
contains $x$. Since $x \leq y$, we have $y \in f^{-1}(V)$, so $f(y) \in V$. Thus
$f(x) \leq_Y f(y)$.
\end{proof}

\begin{warning}
The converse is false in general: a monotone function between $T_0$-spaces need not be
continuous. Continuity is a \emph{stronger} condition than monotonicity. However, for
Alexandrov topologies, the two notions coincide.
\end{warning}

% ------------------------------------------------------------
\section{The Specialization Functor}
% ------------------------------------------------------------

\begin{definition}[Specialization functor]
The \emph{specialization functor} $\Sigma\colon \mathbf{Top}_0 \to \mathbf{Pos}$ assigns
to each $T_0$-space its underlying set with the specialization order, and to each
continuous map the same function (which is monotone by the preceding theorem).
\end{definition}

\begin{theorem}[Left adjoint of the specialization functor]
The Alexandrov functor $\mathrm{Alex}\colon \mathbf{Pos} \to \mathbf{Top}_0$ is left
adjoint to the specialization functor $\Sigma$:
\[
  \mathrm{Alex} \dashv \Sigma.
\]
That is, for every poset $P$ and every $T_0$-space $X$:
\[
  \mathrm{Hom}_{\mathbf{Top}_0}(\mathrm{Alex}(P), X) \cong
  \mathrm{Hom}_{\mathbf{Pos}}(P, \Sigma(X)).
\]
\end{theorem}

\begin{proof}
A morphism $\mathrm{Alex}(P) \to X$ in $\mathbf{Top}_0$ is a continuous map, hence
monotone from $(P,\leq)$ to $(\Sigma(X), \leq_X)$. Conversely, a monotone map
$f\colon P \to \Sigma(X)$ is continuous $\mathrm{Alex}(P) \to X$ because the preimage of
an open $U$ of $X$ (upward-closed for $\leq_X$) is upward-closed in $P$ (since $f$ is
monotone), hence Alexandrov-open.
\end{proof}

% ------------------------------------------------------------
\section{$d$-Spaces}
% ------------------------------------------------------------

\begin{definition}[$d$-space]
A topological space $X$ is a \emph{$d$-space} (or monotone convergence space) if for every
directed set $D$ in $(X, \leq)$ (specialization order) having a supremum $\bigvee D$ in
$(X,\leq)$, every open set containing $\bigvee D$ contains an element of $D$.
Equivalently, if the net along $D$ converges to $\bigvee D$ in the topological sense.
\end{definition}

\begin{theorem}
Every sober space is a $d$-space. Every $T_1$-space is a $d$-space (trivially, since
directed sets for the discrete order are singletons).
\end{theorem}

\begin{proof}
Let $X$ be sober and $D$ a directed set with supremum $s = \bigvee D$. The closure
$\overline{D}$ is a closed set, and it is downward-closed. We show $\overline{D}$ is
irreducible. If $\overline{D} = F_1 \cup F_2$ with $F_1, F_2$ closed, then
$D \subseteq F_1 \cup F_2$. Since $D$ is directed and the $F_i$ are downward-closed, if
$D \not\subseteq F_1$ and $D \not\subseteq F_2$, pick $d_1 \in D \setminus F_1$ and
$d_2 \in D \setminus F_2$. By directedness, there exists $d \geq d_1, d_2$ in $D$. Then
$d \in F_1 \cup F_2$, say $d \in F_1$. Since $F_1$ is downward-closed and $d_1 \leq d$,
we get $d_1 \in F_1$, contradiction. So $\overline{D}$ is irreducible. By sobriety,
$\overline{D}$ has a generic point, which must be $s$.
\end{proof}

% ------------------------------------------------------------
\section{Upper Topology}
% ------------------------------------------------------------

\begin{definition}[Upper topology]
Let $(P, \leq)$ be a partially ordered set. The \emph{upper topology} is the topology
generated by the complements of principal downward-closed sets:
\[
  \tau_{\mathrm{up}} = \langle P \setminus \downarrow\! x : x \in P \rangle.
\]
This is in general coarser than the Alexandrov topology.
\end{definition}

\begin{proposition}
The upper topology has the same specialization order as the Alexandrov topology, namely
the given order $\leq$. It is the coarsest topology with this specialization order.
\end{proposition}

\begin{proof}
Opens of $\tau_{\mathrm{up}}$ are upward-closed, so the specialization order refines
$\leq$. Conversely, if $x \leq y$, then every open $P \setminus \downarrow\! z$ containing
$x$ (i.e., $x \not\leq z$) satisfies $y \not\leq z$ (otherwise $x \leq y \leq z$,
contradiction), so $y$ is in this open. Thus the specialization order is exactly $\leq$.
\end{proof}

% ------------------------------------------------------------
\section{Saturated Sets}
% ------------------------------------------------------------

\begin{definition}[Saturated set]
A subset $A$ of a topological space is \emph{saturated} if it equals the intersection of
all open sets containing it:
$A = \bigcap\{U \in \tau : A \subseteq U\}$.
\end{definition}

\begin{proposition}
In any topological space:
\begin{enumerate}
  \item A subset is saturated if and only if it is upward-closed for the specialization
    order.
  \item Every open set is saturated.
  \item The saturation of a set $A$ is $\uparrow\! A$.
  \item A compact saturated set is the correct analogue of a compact set in the
    non-Hausdorff setting.
\end{enumerate}
\end{proposition}

\begin{proof}
(1) The intersection of all opens containing $A$ is the set of points $y$ such that every
open containing $A$ contains $y$. Now $y$ belongs to every open containing $x \in A$ iff
$x \leq y$. So this intersection is $\uparrow\! A$.
\end{proof}

% ------------------------------------------------------------
\section{Stably Compact Spaces (Preview)}
% ------------------------------------------------------------

\begin{definition}[Stably compact space]
A $T_0$-space is \emph{stably compact} if it is:
\begin{enumerate}
  \item compact,
  \item locally compact (every point has a base of compact saturated neighborhoods),
  \item sober,
  \item \emph{coherent}: the intersection of two compact saturated sets is compact
    saturated.
\end{enumerate}
\end{definition}

\begin{remark}
Stably compact spaces play a central role in the theory, as they are precisely the compact
$T_0$-spaces whose lattice of opens is a \emph{continuous lattice}. We will return to them
in Chapter~8.
\end{remark}

% ------------------------------------------------------------
\section{Exercises}
% ------------------------------------------------------------

\begin{exercise}
Show that a space is $T_1$ if and only if its specialization order is equality.
\end{exercise}

\begin{exercise}
Let $X = \mathrm{Spec}(A)$ be the spectrum of a commutative ring with the Zariski
topology. Describe the specialization order. Show that the generic point of $X$
(if it exists) corresponds to the minimal prime ideal.
\end{exercise}

\begin{exercise}
Show that in a $T_0$-space, the closure of a singleton $\overline{\{x\}}$ is always an
irreducible closed set.
\end{exercise}

\begin{exercise}
Construct a $T_0$-space containing an irreducible closed set that is not the closure of
any singleton. (Hint: consider an ordered set with no greatest element, endowed with the
upper topology.)
\end{exercise}

\begin{exercise}
Let $f\colon X \to Y$ be a continuous open map between $T_0$-spaces. Show that $f$ is
monotone and that $f^{-1}(\uparrow\! y) = \uparrow\! f^{-1}(y)$ for every $y \in Y$.
\end{exercise}

\begin{exercise}[Alexandrov topology and direct limits]
Let $(P_i, f_{ij})_{i \leq j}$ be a direct system of finite posets with monotone maps.
Show that the direct limit in $\mathbf{Top}$ of the corresponding Alexandrov spaces is a
space whose specialization order is the direct limit of the orders.
\end{exercise}
